\nuevaReceta{Fritaílla Almeriense}{60 min}{receta:fritailla}

    \begin{minipage}[t]{0.35\textwidth}
        \section*{Ingredientes}
        \begin{itemize}[leftmargin=*]
            \item 1 kg de tomates maduros
            \item 2 Pimientos verdes
            \item 1 Pimiento rojo
            \item 1 Cebolla
            \item 1 Calabacín
            \item 1 Berenjena
            \item 2 Dientes de ajo
            \item Aceite de oliva
            \item Sal
            \item Una pizca de azúcar (opcional)
        \end{itemize}
    \end{minipage}
    \hfill
    \begin{minipage}[t]{0.6\textwidth}
        \section*{Preparación}
        \begin{enumerate}
            \item \textbf{Preparar las verduras:} Pelar la cebolla y lavar el resto de las verduras. Cortarlo todo en trozos pequeños y regulares.
            \item \textbf{Preparar la berenjena:} Salar ligeramente la berenjena y dejarla reposar 20 minutos. Enjuagarla y secarla para reducir su amargor.
            \item \textbf{Empezar el sofrito:} Calentar una cantidad generosa de aceite en una sartén amplia. Añadir los ajos y la cebolla, y cocinar a fuego medio hasta que estén tiernos.
            \item \textbf{Añadir los pimientos:} Incorporar los pimientos y cocinar durante 10 minutos, removiendo ocasionalmente.
            \item \textbf{Incorporar el resto:} Añadir la berenjena y el calabacín. Continuar la cocción hasta que todas las verduras estén blandas.
            \item \textbf{Añadir el tomate:} Incorporar los tomates pelados y troceados o rallados. Salar y, si están muy ácidos, añadir una pizca de azúcar.
            \item \textbf{Reducir:} Cocinar a fuego lento durante 25-30 minutos, removiendo de vez en cuando, hasta que el tomate haya reducido y el aceite vuelva a aparecer en la superficie.
            \item \textbf{Reposar:} Ajustar de sal y dejar reposar unos minutos antes de servir caliente o a temperatura ambiente.
        \end{enumerate}
    \end{minipage}

    \vspace{1em}
    \begin{tcolorbox}[colback=orange!5, colframe=orange!50, title=Consejo, size=small]
        La fritaílla mejora después de unas horas de reposo. Puede servirse sola con pan, como guarnición o acompañada de huevo, carne o pescado.
    \end{tcolorbox}
\end{receta}
